% Protocol account grounded in the companion executable and formal definitions.
% Bundled paper template assets and bibliography.
\begin{filecontents*}{arXiv-template-LICENSE.txt}
                                 Apache License
                           Version 2.0, January 2004
                        http://www.apache.org/licenses/

   TERMS AND CONDITIONS FOR USE, REPRODUCTION, AND DISTRIBUTION

   1. Definitions.

      "License" shall mean the terms and conditions for use, reproduction,
      and distribution as defined by Sections 1 through 9 of this document.

      "Licensor" shall mean the copyright owner or entity authorized by
      the copyright owner that is granting the License.

      "Legal Entity" shall mean the union of the acting entity and all
      other entities that control, are controlled by, or are under common
      control with that entity. For the purposes of this definition,
      "control" means (i) the power, direct or indirect, to cause the
      direction or management of such entity, whether by contract or
      otherwise, or (ii) ownership of fifty percent (50%) or more of the
      outstanding shares, or (iii) beneficial ownership of such entity.

      "You" (or "Your") shall mean an individual or Legal Entity
      exercising permissions granted by this License.

      "Source" form shall mean the preferred form for making modifications,
      including but not limited to software source code, documentation
      source, and configuration files.

      "Object" form shall mean any form resulting from mechanical
      transformation or translation of a Source form, including but
      not limited to compiled object code, generated documentation,
      and conversions to other media types.

      "Work" shall mean the work of authorship, whether in Source or
      Object form, made available under the License, as indicated by a
      copyright notice that is included in or attached to the work
      (an example is provided in the Appendix below).

      "Derivative Works" shall mean any work, whether in Source or Object
      form, that is based on (or derived from) the Work and for which the
      editorial revisions, annotations, elaborations, or other modifications
      represent, as a whole, an original work of authorship. For the purposes
      of this License, Derivative Works shall not include works that remain
      separable from, or merely link (or bind by name) to the interfaces of,
      the Work and Derivative Works thereof.

      "Contribution" shall mean any work of authorship, including
      the original version of the Work and any modifications or additions
      to that Work or Derivative Works thereof, that is intentionally
      submitted to Licensor for inclusion in the Work by the copyright owner
      or by an individual or Legal Entity authorized to submit on behalf of
      the copyright owner. For the purposes of this definition, "submitted"
      means any form of electronic, verbal, or written communication sent
      to the Licensor or its representatives, including but not limited to
      communication on electronic mailing lists, source code control systems,
      and issue tracking systems that are managed by, or on behalf of, the
      Licensor for the purpose of discussing and improving the Work, but
      excluding communication that is conspicuously marked or otherwise
      designated in writing by the copyright owner as "Not a Contribution."

      "Contributor" shall mean Licensor and any individual or Legal Entity
      on behalf of whom a Contribution has been received by Licensor and
      subsequently incorporated within the Work.

   2. Grant of Copyright License. Subject to the terms and conditions of
      this License, each Contributor hereby grants to You a perpetual,
      worldwide, non-exclusive, no-charge, royalty-free, irrevocable
      copyright license to reproduce, prepare Derivative Works of,
      publicly display, publicly perform, sublicense, and distribute the
      Work and such Derivative Works in Source or Object form.

   3. Grant of Patent License. Subject to the terms and conditions of
      this License, each Contributor hereby grants to You a perpetual,
      worldwide, non-exclusive, no-charge, royalty-free, irrevocable
      (except as stated in this section) patent license to make, have made,
      use, offer to sell, sell, import, and otherwise transfer the Work,
      where such license applies only to those patent claims licensable
      by such Contributor that are necessarily infringed by their
      Contribution(s) alone or by combination of their Contribution(s)
      with the Work to which such Contribution(s) was submitted. If You
      institute patent litigation against any entity (including a
      cross-claim or counterclaim in a lawsuit) alleging that the Work
      or a Contribution incorporated within the Work constitutes direct
      or contributory patent infringement, then any patent licenses
      granted to You under this License for that Work shall terminate
      as of the date such litigation is filed.

   4. Redistribution. You may reproduce and distribute copies of the
      Work or Derivative Works thereof in any medium, with or without
      modifications, and in Source or Object form, provided that You
      meet the following conditions:

      (a) You must give any other recipients of the Work or
          Derivative Works a copy of this License; and

      (b) You must cause any modified files to carry prominent notices
          stating that You changed the files; and

      (c) You must retain, in the Source form of any Derivative Works
          that You distribute, all copyright, patent, trademark, and
          attribution notices from the Source form of the Work,
          excluding those notices that do not pertain to any part of
          the Derivative Works; and

      (d) If the Work includes a "NOTICE" text file as part of its
          distribution, then any Derivative Works that You distribute must
          include a readable copy of the attribution notices contained
          within such NOTICE file, excluding those notices that do not
          pertain to any part of the Derivative Works, in at least one
          of the following places: within a NOTICE text file distributed
          as part of the Derivative Works; within the Source form or
          documentation, if provided along with the Derivative Works; or,
          within a display generated by the Derivative Works, if and
          wherever such third-party notices normally appear. The contents
          of the NOTICE file are for informational purposes only and
          do not modify the License. You may add Your own attribution
          notices within Derivative Works that You distribute, alongside
          or as an addendum to the NOTICE text from the Work, provided
          that such additional attribution notices cannot be construed
          as modifying the License.

      You may add Your own copyright statement to Your modifications and
      may provide additional or different license terms and conditions
      for use, reproduction, or distribution of Your modifications, or
      for any such Derivative Works as a whole, provided Your use,
      reproduction, and distribution of the Work otherwise complies with
      the conditions stated in this License.

   5. Submission of Contributions. Unless You explicitly state otherwise,
      any Contribution intentionally submitted for inclusion in the Work
      by You to the Licensor shall be under the terms and conditions of
      this License, without any additional terms or conditions.
      Notwithstanding the above, nothing herein shall supersede or modify
      the terms of any separate license agreement you may have executed
      with Licensor regarding such Contributions.

   6. Trademarks. This License does not grant permission to use the trade
      names, trademarks, service marks, or product names of the Licensor,
      except as required for reasonable and customary use in describing the
      origin of the Work and reproducing the content of the NOTICE file.

   7. Disclaimer of Warranty. Unless required by applicable law or
      agreed to in writing, Licensor provides the Work (and each
      Contributor provides its Contributions) on an "AS IS" BASIS,
      WITHOUT WARRANTIES OR CONDITIONS OF ANY KIND, either express or
      implied, including, without limitation, any warranties or conditions
      of TITLE, NON-INFRINGEMENT, MERCHANTABILITY, or FITNESS FOR A
      PARTICULAR PURPOSE. You are solely responsible for determining the
      appropriateness of using or redistributing the Work and assume any
      risks associated with Your exercise of permissions under this License.

   8. Limitation of Liability. In no event and under no legal theory,
      whether in tort (including negligence), contract, or otherwise,
      unless required by applicable law (such as deliberate and grossly
      negligent acts) or agreed to in writing, shall any Contributor be
      liable to You for damages, including any direct, indirect, special,
      incidental, or consequential damages of any character arising as a
      result of this License or out of the use or inability to use the
      Work (including but not limited to damages for loss of goodwill,
      work stoppage, computer failure or malfunction, or any and all
      other commercial damages or losses), even if such Contributor
      has been advised of the possibility of such damages.

   9. Accepting Warranty or Additional Liability. While redistributing
      the Work or Derivative Works thereof, You may choose to offer,
      and charge a fee for, acceptance of support, warranty, indemnity,
      or other liability obligations and/or rights consistent with this
      License. However, in accepting such obligations, You may act only
      on Your own behalf and on Your sole responsibility, not on behalf
      of any other Contributor, and only if You agree to indemnify,
      defend, and hold each Contributor harmless for any liability
      incurred by, or claims asserted against, such Contributor by reason
      of your accepting any such warranty or additional liability.

   END OF TERMS AND CONDITIONS

   APPENDIX: How to apply the Apache License to your work.

      To apply the Apache License to your work, attach the following
      boilerplate notice, with the fields enclosed by brackets "[]"
      replaced with your own identifying information. (Don't include
      the brackets!)  The text should be enclosed in the appropriate
      comment syntax for the file format. We also recommend that a
      file or class name and description of purpose be included on the
      same "printed page" as the copyright notice for easier
      identification within third-party archives.

   Copyright [yyyy] [name of copyright owner]

   Licensed under the Apache License, Version 2.0 (the "License");
   you may not use this file except in compliance with the License.
   You may obtain a copy of the License at

       http://www.apache.org/licenses/LICENSE-2.0

   Unless required by applicable law or agreed to in writing, software
   distributed under the License is distributed on an "AS IS" BASIS,
   WITHOUT WARRANTIES OR CONDITIONS OF ANY KIND, either express or implied.
   See the License for the specific language governing permissions and
   limitations under the License.
\end{filecontents*}

\begin{filecontents*}{preprint.sty}
\NeedsTeXFormat{LaTeX2e}

%% declare nonatbib option, which does not load natbib in case of
%% package clash (users can pass options to natbib via
%% \PassOptionsToPackage)
%\newif\if@natbib\@natbibtrue
%\DeclareOption{nonatbib}{
%  \@natbibfalse
%}

\ProcessOptions\relax

% Use roman fonts for equations if available
\IfFileExists{txfonts.sty}%
  {\AtEndOfClass{\RequirePackage{txfonts}%
  \gdef\ttdefault{cmtt}%
   \let\iint\relax
  \let\iiint\relax
  \let\iiiint\relax
  \let\idotsint\relax
  \let\openbox\relax}}{\RequirePackage{times}}

 % other font configurations
\renewcommand{\rmdefault}{ptm}
\renewcommand{\sfdefault}{phv}


\RequirePackage[verbose=true,letterpaper]{geometry}
\AtBeginDocument{
  \newgeometry{
    textheight=9.3in,
    textwidth=7.1in,
    top=0.96in,
    headheight=0.18in,
    headsep=0.14in,
    footskip=0.1in
  }
}

\usepackage{tikz}
\usepackage{eso-pic}

\AddToHook{shipout/background}{%
  \begin{tikzpicture}[remember picture, overlay]
    \node at ([xshift=-3.55in,yshift=5.2in]current page.center) {\rule{.4pt}{.4in}};
    \node at ([xshift=3.55in,yshift=5.2in]current page.center) {\rule{.4pt}{.4in}};
    \node at ([xshift=-3.95in,yshift=4.8in]current page.center) {\rule{.4in}{.4pt}};
    \node at ([xshift=3.95in,yshift=4.8in]current page.center) {\rule{.4in}{.4pt}};
    \node at ([xshift=-3.55in,yshift=-5.2in]current page.center) {\rule{.4pt}{.4in}};
    \node at ([xshift=3.55in,yshift=-5.2in]current page.center) {\rule{.4pt}{.4in}};
    \node at ([xshift=-3.95in,yshift=-4.8in]current page.center) {\rule{.4in}{.4pt}};
    \node at ([xshift=3.95in,yshift=-4.8in]current page.center) {\rule{.4in}{.4pt}};
  \end{tikzpicture}%
}


\widowpenalty=10000
\clubpenalty=10000
\flushbottom
\sloppy

% easy adding of ORCiDs
\usepackage{scalerel}
\usepackage{tikz}
\usetikzlibrary{svg.path}
\definecolor{orcidlogocol}{HTML}{A6CE39}
\tikzset{
  orcidlogo/.pic={
    \fill[orcidlogocol] svg{M256,128c0,70.7-57.3,128-128,128C57.3,256,0,198.7,0,128C0,57.3,57.3,0,128,0C198.7,0,256,57.3,256,128z};
    \fill[white] svg{M86.3,186.2H70.9V79.1h15.4v48.4V186.2z}
                 svg{M108.9,79.1h41.6c39.6,0,57,28.3,57,53.6c0,27.5-21.5,53.6-56.8,53.6h-41.8V79.1z M124.3,172.4h24.5c34.9,0,42.9-26.5,42.9-39.7c0-21.5-13.7-39.7-43.7-39.7h-23.7V172.4z}
                 svg{M88.7,56.8c0,5.5-4.5,10.1-10.1,10.1c-5.6,0-10.1-4.6-10.1-10.1c0-5.6,4.5-10.1,10.1-10.1C84.2,46.7,88.7,51.3,88.7,56.8z};
  }
}

\newcommand\orcid[1]{\href{https://orcid.org/#1}{\mbox{\scalerel*{
\begin{tikzpicture}[yscale=-1,transform shape]
\pic{orcidlogo};
\end{tikzpicture}
}{|}}} \href{https://orcid.org/#1}{#1}}

% Header options
\RequirePackage{fancyhdr}
\fancyhf{}
\pagestyle{fancy}
\renewcommand{\headrulewidth}{0pt}
\fancyheadoffset{0pt}
\lhead{\scshape Preprint -- \@title}
\rhead{\thepage}
%\rfoot{\thepage}

%Handling Keywords
\def\keywordname{{\bfseries \emph Keywords}}%
\def\keywords#1{\par\addvspace\medskipamount{\rightskip=0pt plus1cm
\def\and{\ifhmode\unskip\nobreak\fi\ $\cdot$
}\noindent\keywordname\enspace\ignorespaces#1\par}}

% font sizes with reduced leading
\renewcommand{\normalsize}{%
  \@setfontsize\normalsize\@xpt\@xipt
  \abovedisplayskip      7\p@ \@plus 2\p@ \@minus 5\p@
  \abovedisplayshortskip \z@ \@plus 3\p@
  \belowdisplayskip      \abovedisplayskip
  \belowdisplayshortskip 4\p@ \@plus 3\p@ \@minus 3\p@
}
\normalsize
\renewcommand{\small}{%
  \@setfontsize\small\@ixpt\@xpt
  \abovedisplayskip      6\p@ \@plus 1.5\p@ \@minus 4\p@
  \abovedisplayshortskip \z@  \@plus 2\p@
  \belowdisplayskip      \abovedisplayskip
  \belowdisplayshortskip 3\p@ \@plus 2\p@   \@minus 2\p@
}
\renewcommand{\footnotesize}{\@setfontsize\footnotesize\@ixpt\@xpt}
\renewcommand{\scriptsize}{\@setfontsize\scriptsize\@viipt\@viiipt}
\renewcommand{\tiny}{\@setfontsize\tiny\@vipt\@viipt}
\renewcommand{\large}{\@setfontsize\large\@xiipt{14}}
\renewcommand{\Large}{\@setfontsize\Large\@xivpt{16}}
\renewcommand{\LARGE}{\@setfontsize\LARGE\@xviipt{20}}
\renewcommand{\huge}{\@setfontsize\huge\@xxpt{23}}
\renewcommand{\Huge}{\@setfontsize\Huge\@xxvpt{28}}

% sections with less space
\providecommand{\section}{}
\renewcommand{\section}{%
  \@startsection{section}{1}{\z@}%
                {-2.0ex \@plus -0.5ex \@minus -0.2ex}%
                { 1.5ex \@plus  0.3ex \@minus  0.2ex}%
                {\large\scshape\raggedright}%
}
\providecommand{\subsection}{}
\renewcommand{\subsection}{%
  \@startsection{subsection}{2}{\z@}%
                {-1.8ex \@plus -0.5ex \@minus -0.2ex}%
                { 0.8ex \@plus  0.2ex}%
                {\normalsize\itshape\raggedright}%
}
\providecommand{\subsubsection}{}
\renewcommand{\subsubsection}{%
  \@startsection{subsubsection}{3}{\z@}%
                {-1.5ex \@plus -0.5ex \@minus -0.2ex}%
                { 0.5ex \@plus  0.2ex}%
                {\normalsize\itshape\raggedright}%
}
\providecommand{\paragraph}{}
\renewcommand{\paragraph}{%
  \@startsection{paragraph}{4}{\z@}%
                {1.5ex \@plus 0.5ex \@minus 0.2ex}%
                {-1em}%
                {\normalsize\bf}%
}
\providecommand{\subparagraph}{}
\renewcommand{\subparagraph}{%
  \@startsection{subparagraph}{5}{\z@}%
                {1.5ex \@plus 0.5ex \@minus 0.2ex}%
                {-1em}%
                {\normalsize\bf}%
}
\providecommand{\subsubsubsection}{}
\renewcommand{\subsubsubsection}{%
  \vskip5pt{\noindent\normalsize\rm\raggedright}%
}

% float placement
\renewcommand{\topfraction      }{0.85}
\renewcommand{\bottomfraction   }{0.4}
\renewcommand{\textfraction     }{0.1}
\renewcommand{\floatpagefraction}{0.7}

\newlength{\@abovecaptionskip}\setlength{\@abovecaptionskip}{7\p@}
\newlength{\@belowcaptionskip}\setlength{\@belowcaptionskip}{\z@}

\setlength{\abovecaptionskip}{\@abovecaptionskip}
\setlength{\belowcaptionskip}{\@belowcaptionskip}

% swap above/belowcaptionskip lengths for tables
\renewenvironment{table}
  {\setlength{\abovecaptionskip}{\@belowcaptionskip}%
   \setlength{\belowcaptionskip}{\@abovecaptionskip}%
   \@float{table}}
  {\end@float}

% footnote formatting
\setlength{\footnotesep }{6.65\p@}
\setlength{\skip\footins}{9\p@ \@plus 4\p@ \@minus 2\p@}
\renewcommand{\footnoterule}{\kern-3\p@ \hrule width 12pc \kern 2.6\p@}
\setcounter{footnote}{0}

% paragraph formatting
\setlength{\parindent}{2em}
\setlength{\parskip  }{5.5\p@}

% list formatting
\setlength{\topsep       }{4\p@ \@plus 1\p@   \@minus 2\p@}
\setlength{\partopsep    }{1\p@ \@plus 0.5\p@ \@minus 0.5\p@}
\setlength{\itemsep      }{2\p@ \@plus 1\p@   \@minus 0.5\p@}
\setlength{\parsep       }{2\p@ \@plus 1\p@   \@minus 0.5\p@}
\setlength{\leftmargin   }{3pc}
\setlength{\leftmargini  }{\leftmargin}
\setlength{\leftmarginii }{2em}
\setlength{\leftmarginiii}{1.5em}
\setlength{\leftmarginiv }{1.0em}
\setlength{\leftmarginv  }{0.5em}
\def\@listi  {\leftmargin\leftmargini}
\def\@listii {\leftmargin\leftmarginii
              \labelwidth\leftmarginii
              \advance\labelwidth-\labelsep
              \topsep  2\p@ \@plus 1\p@    \@minus 0.5\p@
              \parsep  1\p@ \@plus 0.5\p@ \@minus 0.5\p@
              \itemsep \parsep}
\def\@listiii{\leftmargin\leftmarginiii
              \labelwidth\leftmarginiii
              \advance\labelwidth-\labelsep
              \topsep    1\p@ \@plus 0.5\p@ \@minus 0.5\p@
              \parsep    \z@
              \partopsep 0.5\p@ \@plus 0\p@ \@minus 0.5\p@
              \itemsep \topsep}
\def\@listiv {\leftmargin\leftmarginiv
              \labelwidth\leftmarginiv
              \advance\labelwidth-\labelsep}
\def\@listv  {\leftmargin\leftmarginv
              \labelwidth\leftmarginv
              \advance\labelwidth-\labelsep}
\def\@listvi {\leftmargin\leftmarginvi
              \labelwidth\leftmarginvi
              \advance\labelwidth-\labelsep}

% create title
\providecommand{\maketitle}{}
\renewcommand{\maketitle}{%
  \par
  \begingroup
    \renewcommand{\thefootnote}{\fnsymbol{footnote}}
    % for perfect author name centering
    \renewcommand{\@makefnmark}{\hbox to \z@{$^{\@thefnmark}$\hss}}
    % The footnote-mark was overlapping the footnote-text,
    % added the following to fix this problem               (MK)
    \long\def\@makefntext##1{%
      \parindent 1em\noindent
      \hbox to 1.8em{\hss $\m@th ^{\@thefnmark}$}##1
    }
    \thispagestyle{empty}
    \vspace*{-0.5cm}
    \@maketitle
    \@thanks
    %\@notice
  \endgroup
  \let\maketitle\relax
  \let\thanks\relax
}

% create title (includes both anonymized and non-anonymized versions)
\providecommand{\@maketitle}{}
\renewcommand{\@maketitle}{%
  \vbox{%
    \hsize\textwidth
    \linewidth\hsize
    \centering
    {\LARGE\sc \@title\par}
    \vskip 0.1in
    \textsc{Preprint, compiled \today}\\
    \def\And{%
      \end{tabular}\hfil\linebreak[0]\hfil%
      \begin{tabular}[t]{c}\bf\rule{\z@}{24\p@}\ignorespaces%
    }
    \def\AND{%
      \end{tabular}\hfil\linebreak[4]\hfil%
      \begin{tabular}[t]{c}\bf\rule{\z@}{24\p@}\ignorespaces%
    }
    \begin{tabular}[t]{c}\bf\rule{\z@}{24\p@}\@author\end{tabular}%
  \vskip 0.2in
  }
}

% abstract styling
\renewenvironment{abstract}
{
  \centerline
  {\large \bfseries \scshape Abstract}
  \begin{quote}
}
{
  \end{quote}
}

\endinput
\end{filecontents*}

\begin{filecontents*}{references.bib}
@techreport{rw-vrf,
  author = {Sharon Goldberg and Leonid Reyzin and Dimitrios Papadopoulos and Jan V{\v c}el{\'a}k},
  title = {Verifiable Random Functions ({VRFs})},
  institution = {Internet Research Task Force},
  type = {RFC},
  number = {9381},
  year = {2023},
  month = aug,
  url = {https://www.rfc-editor.org/rfc/rfc9381.html}
}

@article{rw-dls,
  author = {Cynthia Dwork and Nancy Lynch and Larry Stockmeyer},
  title = {Consensus in the Presence of Partial Synchrony},
  journal = {Journal of the ACM},
  volume = {35},
  number = {2},
  year = {1988},
  pages = {288--323},
  url = {https://groups.csail.mit.edu/tds/papers/Lynch/jacm88.pdf}
}

@inproceedings{rw-pbft,
  author = {Miguel Castro and Barbara Liskov},
  title = {Practical Byzantine Fault Tolerance},
  booktitle = {Proceedings of the 3rd Symposium on Operating Systems Design and Implementation},
  year = {1999},
  pages = {173--186},
  publisher = {USENIX Association},
  doi = {10.5555/296806.296824}
}

@article{rw-tendermint,
  author = {Ethan Buchman and Jae Kwon and Zarko Milosevic},
  title = {The Latest Gossip on {BFT} Consensus},
  journal = {arXiv preprint arXiv:1807.04938},
  year = {2019},
  doi = {10.48550/arXiv.1807.04938}
}

@inproceedings{rw-hotstuff,
  author = {Maofan Yin and Dahlia Malkhi and Michael K. Reiter and Guy Golan Gueta and Ittai Abraham},
  title = {{HotStuff}: {BFT} Consensus in the Lens of Blockchain},
  booktitle = {Proceedings of the 2019 ACM Symposium on Principles of Distributed Computing},
  year = {2019},
  pages = {347--356},
  doi = {10.1145/3293611.3331591},
  url = {https://arxiv.org/pdf/1803.05069v6}
}

@techreport{rw-diembft,
  author = {{The Diem Team}},
  title = {{DiemBFT} v4: State Machine Replication in the {Diem} Blockchain},
  institution = {Diem Association},
  year = {2021},
  month = aug,
  note = {Revised August 18, 2021},
  url = {https://developers.diem.com/papers/diem-consensus-state-machine-replication-in-the-diem-blockchain/2021-08-17.pdf}
}

@article{rw-monadbft,
  author = {Mohammad Mussadiq Jalalzai and Kushal Babel and Jovan Komatovic and Tobias Klenze and Sourav Das and Fatima Elsheimy and Mike Setrin and John Bergschneider and Babak Gilkalaye},
  title = {{MonadBFT}: Fast, Responsive, Fork-Resistant Streamlined Consensus},
  journal = {arXiv preprint arXiv:2502.20692},
  year = {2026},
  doi = {10.48550/arXiv.2502.20692}
}

@article{rw-casper,
  author = {Vitalik Buterin and Virgil Griffith},
  title = {Casper the Friendly Finality Gadget},
  journal = {arXiv preprint arXiv:1710.09437},
  year = {2019},
  doi = {10.48550/arXiv.1710.09437}
}

@article{rw-gasper,
  author = {Vitalik Buterin and Diego Hernandez and Thor Kamphefner and Khiem Pham and Zhi Qiao and Danny Ryan and Juhyeok Sin and Ying Wang and Yan X Zhang},
  title = {Combining {GHOST} and Casper},
  journal = {arXiv preprint arXiv:2003.03052},
  year = {2020},
  doi = {10.48550/arXiv.2003.03052}
}

@article{rw-grandpa,
  author = {Alistair Stewart and Eleftherios Kokoris-Kogias},
  title = {{GRANDPA}: A Byzantine Finality Gadget},
  journal = {arXiv preprint arXiv:2007.01560},
  year = {2020},
  doi = {10.48550/arXiv.2007.01560}
}

@inproceedings{rw-backbone,
  author = {Juan A. Garay and Aggelos Kiayias and Nikos Leonardos},
  title = {The Bitcoin Backbone Protocol: Analysis and Applications},
  booktitle = {Advances in Cryptology -- EUROCRYPT 2015},
  year = {2015},
  pages = {281--310},
  publisher = {Springer},
  doi = {10.1007/978-3-662-46803-6_10}
}

@inproceedings{rw-narwhal-tusk,
  author = {George Danezis and Lefteris Kokoris-Kogias and Alberto Sonnino and Alexander Spiegelman},
  title = {Narwhal and Tusk: A {DAG}-based Mempool and Efficient {BFT} Consensus},
  booktitle = {Proceedings of the Seventeenth European Conference on Computer Systems},
  year = {2022},
  pages = {1--17},
  doi = {10.1145/3492321.3519594}
}

@article{rw-aleph,
  author = {Adam Gągol and Damian Leśniak and Damian Straszak and Michał Świętek},
  title = {Aleph: Efficient Atomic Broadcast in Asynchronous Networks with Byzantine Nodes},
  journal = {arXiv preprint arXiv:1908.05156},
  year = {2019},
  doi = {10.48550/arXiv.1908.05156}
}

@article{rw-hotstuff2,
  author = {Dahlia Malkhi and Kartik Nayak},
  title = {{HotStuff-2}: Optimal Two-Phase Responsive {BFT}},
  journal = {Cryptology ePrint Archive, Paper 2023/397},
  year = {2023},
  url = {https://eprint.iacr.org/2023/397}
}

@inproceedings{rw-sbft,
  author = {Guy Golan Gueta and Ittai Abraham and Shelly Grossman and Dahlia Malkhi and Benny Pinkas and Michael K. Reiter and Dragos-Adrian Seredinschi and Orr Tamir and Alin Tomescu},
  title = {{SBFT}: A Scalable and Decentralized Trust Infrastructure},
  booktitle = {2019 49th Annual IEEE/IFIP International Conference on Dependable Systems and Networks},
  year = {2019},
  pages = {568--580},
  doi = {10.1109/DSN.2019.00063}
}

@inproceedings{rw-simplex,
  author = {Benjamin Y. Chan and Rafael Pass},
  title = {Simplex Consensus: A Simple and Fast Consensus Protocol},
  booktitle = {Theory of Cryptography},
  year = {2023},
  pages = {452--479},
  doi = {10.1007/978-3-031-48624-1_17}
}

@misc{rw-alpenglow,
  author = {Quentin Kniep and Jakub Sliwinski and Roger Wattenhofer},
  title = {Solana Alpenglow Consensus: Increased Bandwidth, Reduced Latency},
  year = {2025},
  howpublished = {White Paper v1.1, Anza},
  url = {https://github.com/rogerANZA/Alpenglow-White-Paper/blob/f29dbbc51af8d83176f1fc06836ba6aa193d731a/Alpenglow-v1.1.pdf}
}

@inproceedings{rw-algorand,
  author = {Yossi Gilad and Rotem Hemo and Silvio Micali and Georgios Vlachos and Nickolai Zeldovich},
  title = {Algorand: Scaling {Byzantine} Agreements for Cryptocurrencies},
  booktitle = {Proceedings of the 26th ACM Symposium on Operating Systems Principles},
  year = {2017},
  pages = {51--68},
  doi = {10.1145/3132747.3132757}
}

@article{rw-minimmit,
  author = {Brendan Kobayashi Chou and Andrew Lewis-Pye and Patrick O'Grady},
  title = {Minimmit: Fast Finality with Even Faster Blocks},
  journal = {arXiv preprint arXiv:2508.10862},
  year = {2026},
  doi = {10.48550/arXiv.2508.10862}
}

@article{rw-chonkybft,
  author = {Bruno Fran{\c{c}}a and Denis Kolegov and Igor Konnov and Grzegorz Prusak},
  title = {{ChonkyBFT}: Consensus Protocol of {ZKsync}},
  journal = {arXiv preprint arXiv:2503.15380},
  year = {2025},
  doi = {10.48550/arXiv.2503.15380}
}

@article{rw-bullshark,
  author = {Alexander Spiegelman and Neil Giridharan and Alberto Sonnino and Lefteris Kokoris-Kogias},
  title = {Bullshark: {DAG} {BFT} Protocols Made Practical},
  journal = {Proceedings of the 2022 ACM SIGSAC Conference on Computer and Communications Security},
  year = {2022},
  pages = {2705--2718},
  doi = {10.1145/3548606.3559361}
}

@inproceedings{rw-jolteon,
  author = {Rati Gelashvili and Lefteris Kokoris-Kogias and Alberto Sonnino and Alexander Spiegelman and Zhuolun Xiang},
  title = {Jolteon and Ditto: Network-Adaptive Efficient Consensus with Asynchronous Fallback},
  booktitle = {Financial Cryptography and Data Security},
  year = {2022},
  pages = {296--315},
  doi = {10.1007/978-3-031-18283-9_14}
}

@article{rw-fast-hotstuff,
  author = {Mohammad M. Jalalzai and Jianyu Niu and Chen Feng and Fangyu Gai},
  title = {Fast-HotStuff: A Fast and Robust {BFT} Protocol for Blockchains},
  journal = {IEEE Transactions on Dependable and Secure Computing},
  year = {2023},
  doi = {10.1109/TDSC.2023.3308848}
}

@inproceedings{rw-mysticeti,
  author = {Kushal Babel and Andrey Chursin and George Danezis and Anastasios Kichidis and Lefteris Kokoris-Kogias and Arun Koshy and Alberto Sonnino and Mingwei Tian},
  title = {Mysticeti: Reaching the Latency Limits with Uncertified {DAGs}},
  booktitle = {Network and Distributed System Security Symposium},
  year = {2025},
  doi = {10.14722/ndss.2025.240929}
}
\end{filecontents*}

\documentclass[twocolumn]{article}
\usepackage{preprint}
\RemoveFromHook{shipout/background}[preprint]
\raggedbottom
\usepackage[utf8]{inputenc}
\usepackage[T1]{fontenc}
\usepackage{courier}
\usepackage{amsmath,amssymb}
\usepackage{microtype}
\usepackage[numbers,square]{natbib}
\bibliographystyle{unsrtnat}
\usepackage{hyperref}
\hypersetup{
  pdftitle={Tsampi BFT: leaderless voting with quorum-ancestry finality},
  pdfauthor={Tim Watts},
  pdfsubject={Byzantine fault-tolerant consensus},
  pdfkeywords={Byzantine fault tolerance, leaderless consensus, linked voting, quorum ancestry},
  hidelinks}
\setlength{\emergencystretch}{3em}
\setcounter{secnumdepth}{3}
\newcommand{\qid}[1]{\operatorname{Q}_{#1}}
\newcommand{\height}{\operatorname{height}}

\title{Tsampi BFT: leaderless voting with quorum-ancestry finality}
\author{Tim Watts\\\href{mailto:tim@readevalprint.com}{tim@readevalprint.com}}
\date{Published August 17, 2026\\Updated November 1, 2026}
\usepackage{etoolbox}
\makeatletter
\patchcmd{\@maketitle}{\textsc{Preprint, compiled \today}}{\@date}{}%
  {\PackageError{tsampi-bft}{Could not replace template date}{Check preprint.sty title formatting.}}
% Keep the template's spacing while using sentence case and bold headings.
\patchcmd{\@maketitle}{\LARGE\sc}{\LARGE\bfseries}{}%
  {\PackageError{tsampi-bft}{Could not replace title small caps}{Check preprint.sty title formatting.}}
\patchcmd{\section}{\scshape}{\normalfont\bfseries}{}%
  {\PackageError{tsampi-bft}{Could not replace section small caps}{Check preprint.sty section formatting.}}
\patchcmd{\subsection}{\itshape}{\normalfont\bfseries}{}%
  {\PackageError{tsampi-bft}{Could not replace subsection styling}{Check preprint.sty subsection formatting.}}
\patchcmd{\subsubsection}{\itshape}{\normalfont\bfseries}{}%
  {\PackageError{tsampi-bft}{Could not replace subsubsection styling}{Check preprint.sty subsubsection formatting.}}
\patchcmd{\abstract}{\scshape}{}{}%
  {\PackageError{tsampi-bft}{Could not replace abstract small caps}{Check preprint.sty abstract formatting.}}
\lhead{Preprint -- \@title}
\makeatother

\begin{document}
\twocolumn[
  \begin{@twocolumnfalse}
    \maketitle
    \thispagestyle{empty}
    \begin{abstract}
% Fixed opening sentence: preserve this wording exactly in future edits.
Tsampi BFT is a leaderless, asynchronous Byzantine fault tolerant state replication protocol.
A Vote endorses one Block and uses PrevVote and PrevSelfVote to record ancestry and signed order.
Blocks contain Votes endorsing their parent Block, and different children can contain the same Votes.
Ancestry, scores and historical quorum queries determine permitted Votes and finalization.
Under fixed authority and a quorum for each parent, three later Blocks finalize a selected Block; a fourth finalizes its child.
The paper states the assumptions for Safety and conditional progress and compares related research.
\end{abstract}
    \vspace{0.35cm}
  \end{@twocolumnfalse}
]

\section{Introduction}\label{sec:overview}

Tsampi BFT is a leaderless, asynchronous Byzantine fault tolerant state replication protocol.
Safety requires compatible finalized lineages across honest nodes and time.
Liveness requires every honest node to eventually finalize a Block above any specified common reference height.
Validators sign Votes, and proposed Blocks contain Votes endorsing their parent Block.
PrevVote follows Block ancestry, while PrevSelfVote records one Validator's signed order within active linked history.
Nodes use ancestry, scores and quorum queries to select Blocks and advance their finalized Blocks.

\section{Votes, Blocks and Validators}\label{sec:objects}

\subsection{Validators and stake}

An authority specifies the Validators, their stake, the chain identity and the endorsement threshold for admitting a child Block.
A Validator has positive stake under that authority.
A node retains and checks Blocks and Votes; it may sign for a Validator or observe without signing.
For authority $s$ with Validator set $V_s$, let $w_s(i)$ denote Validator $i$'s stake and $W_s$ its total stake.
\[
W_s=\sum_{i\in V_s}w_s(i).
\]
Each check uses the authority recorded by the relevant Block, preserving historical weights after later exclusions.

\subsection{Blocks and parent Votes}\label{sec:proposals}

A non-genesis Block identifies one exact parent Block and contains Votes endorsing that parent under its historical authority.
Each signing Validator occurs once, and the Block references the proposing Validator's parent Vote as its proposer Vote.
The proposing Validator signs a separate detached Vote endorsing the completed child Block.
Requiring this Vote first makes each proposer begin with its own Block and restricts later endorsements to permitted improvements.
Different children can contain the same parent Votes, including children proposed by one Validator.
The genesis Block has height zero, empty parent, proposer Vote and inherited Q3 references, and empty Vote and evidence collections.

For $N$ Validators following the improvement rules, consider one child Block each, with a common parent and equal Q1 and Q2.
If initial endorsements were unrestricted, each Validator could endorse all $N$ children from least to most preferred, giving $N^2$ Votes.
Requiring the detached Vote first makes each proposer start with its own child and endorse only better children afterward.
From most to least preferred child, proposers can sign at most $1,2,\ldots,N$ Votes respectively, bounding the total $V$, including detached Votes.
\[
V\leq\sum_{r=1}^{N}r=\frac{N(N+1)}{2}.
\]
Additional Votes after the detached Votes total at most $N(N-1)/2$.
The leading term falls from $N^2$ to $N^2/2$, reducing signing and verification work for these competing children.

A child Block uses its parent's endorsement threshold $n/d$, with $0<n\leq d$, to determine the required stake.
The default $n/d=2/3$ requires strictly more than two thirds of the parent's stake; $n=d$ requires all stake.
Let $U_B$ be the distinct signing Validators of child $B$'s parent Votes and $A_B$ their total stake under parent $P$'s authority $s_P$.
\[
\begin{aligned}
A_B&=\sum_{i\in U_B}w_{s_P}(i),\\
&\begin{cases}
A_Bd>W_{s_P}n,&n<d,\\
A_B=W_{s_P},&n=d.
\end{cases}
\end{aligned}
\]
The proposing Validator must have positive stake under that parent authority.
Verified evidence determines the child's reduced authority as described in Section~\ref{sec:authority}.
Each child also advances its parent's raw Q1 result or remaining quorum stake under Section~\ref{sec:parent-progress}.

\subsection{Votes and references}

A Vote identifies its signing Validator and one exact endorsed Block, written $T(v)$ for Vote $v$.
Authentication binds the Vote's contents to that Validator and checks its historical stake.
Its chain, height and proposing Validator must match the endorsed Block, with Vote height $\height(T(v))+1$.
The Vote also carries PrevVote, PrevSelfVote, score context and evidence required by its signed references.
A node admits at most one distinct Vote per Validator and exact endorsed Block on a chain.

\subsubsection{PrevVote}

Tsampi's core logic model is a VoteChain: a chain of Votes linked by PrevVote.
Each Block references its proposer Vote, which endorses the parent Block.
PrevVote names the inner identity of that proposer Vote, resolved through the exact endorsed Block.
The inner identity fixes the signing Validator, chain, height, proposing Validator and PrevVote; the complete Vote additionally identifies its endorsed Block.
Following PrevVote therefore follows the signed Votes underlying Block ancestry.
A Vote endorsing genesis has an empty PrevVote reference.

\subsubsection{PrevSelfVote}\label{sec:segments}

During live admission, PrevSelfVote names the preceding signed Vote by the same Validator on the same chain.
These references give the Validator's active linked history a total causal order across endorsed Blocks.
Each successor endorses a Block that improves on its predecessor's endorsed Block under Section~\ref{sec:permission}.
The node also checks improvement against the Blocks established by that Vote's referenced history.
The linked history remains active while any endorsed Block equals or descends from the new endorsed Block's Q3.
Retirement permits an empty PrevSelfVote after a complete check, while recorded history retains every signed Vote.

Catch-up establishes Q3 support from authenticated Votes collected across referenced descendant Blocks, counting distinct Validators under an established historical authority.
With strictly more than two thirds honest stake, each quorum includes honest Validators who already validated those descendants and their active histories.
Catch-up relies on that prior validation for finalized history under the assumptions in Section~\ref{sec:assumptions}.
Fresh Votes after catch-up satisfy the active PrevSelfVote and improvement checks.

\section{Block selection and finalization}

\subsection{Ancestry}

A lineage is the exact sequence of parent-linked Blocks from genesis through a specified Block.
This sequence determines whether finalized lineages are compatible.
Write $A\preceq B$ when $A$ is an ancestor of or equal to $B$, and $A\prec B$ for strict ancestry.
Two lineages are compatible when one is a prefix of the other; otherwise they are incomparable.

\subsection{Q1, Q2 and Q3}\label{sec:quorums}

Q1 scans from an admitted Block's parent toward genesis and accumulates distinct signing Validators from carried parent Votes.
A Vote supports its endorsed Block and that Block's ancestors.
At each Block $T$, Q1 counts each accumulated Validator once under $T$'s historical authority.
For authority $s_T$ and counted Validator set $U$, a quorum requires the fixed strict two-thirds inequality below.
\[
3\sum_{i\in U}w_{s_T}(i)>2W_{s_T}.
\]
Q1 returns the Block with the first completed quorum, unless an inherited limiting Block stops the query earlier.
An inherited limit or genesis fallback identifies the returned Block without establishing a completed quorum.

Each child must advance its parent's raw Q1 result or increase the remaining distinct stake.\label{sec:parent-progress}
This ties Block growth to the quorum queries used for finalization.
The parent Block's Q3 height limits the child's Q1 scan; an empty inherited reference permits scanning to genesis.
The full query replaces an absent raw result with genesis.
Equal raw results require increased remaining distinct stake; different nonempty results require the child's result to strictly descend from the parent's.
A nonempty child result also permits progress when the parent's raw result is absent.

Q2 and Q3 repeat Q1 from the Block returned by the preceding query.
Each repetition uses that Block's own lineage, historical authority and inherited limit.
Each repetition starts with zero counted Validators, as expressed by the following definitions.
\[
\qid{2}(B)=\qid{1}(\qid{1}(B)),\qquad
\qid{3}(B)=\qid{1}(\qid{2}(B)).
\]
The specified Block's ancestry and historical authority make each query reproducible as the node's state advances.

\subsection{Finalized Blocks}\label{sec:finalization}\label{sec:stored-finality}

A node's finalized Block records the ancestry that later finalization must preserve.
Compatible Q3 results from admitted Blocks let the node extend that ancestry.
The node retains equal or deeper compatible Q3 results from admitted lineages extending its existing finalized Block.
Its Tip is an admitted Block selected for possible finalization and must equal or descend from the finalized Block.
As finality advances, a compatible Tip may remain; otherwise the node selects the newly finalized Block as Tip.
The finalized Block can therefore advance beyond the Tip's own historical Q3.

Consider a fixed authority and lineage $G,A,B,C,D,E$, with a quorum of parent Votes carried by every child.
At $D$, the queries return the following Blocks.
\[
\qid{1}(D)=C,\quad \qid{2}(D)=B,\quad \qid{3}(D)=A.
\]
Thus a node admitting $D$ can finalize $A$ after three completed Q1 advances.
A fourth advance gives $\qid{3}(E)=B$ and lets that node finalize the first new child of $A$.
These advances are asynchronous completion events and can take arbitrary time.
The example demonstrates progress beyond $A$ under the delivery and fresh generation assumptions of Section~\ref{sec:liveness}.

\subsection{Block comparison}\label{sec:permission}

Tip selection and signing use the same improvement rule for a change of endorsed Block.
Normal selection chooses an admitted eligible Block equal to or descending from the node's finalized Block.
Signing also preserves the obligations in the Validator's recorded Votes and the order declared through PrevSelfVote.
Verified evidence determines the child's authority, and the improvement rule determines permitted Tip changes.

\subsubsection{Scores}\label{sec:ordering}

The proposer score is a fixed input attached to the proposer Vote and reused by every Block referencing that Vote.
On one lineage, prefer the strict descendant.
On incomparable lineages, compare proposer Votes at the first different child Blocks and prefer the higher score.
For equal scores, prefer the smaller Block identity; later descendants preserve this result.
Write $R(A,B)$ when the comparison prefers Block $A$ to Block $B$.

\subsubsection{Improvement}

Improvement determines whether new Block $N$ may replace old Block $O$ as a Tip or endorsed Block.
On one lineage, improvement requires $O\prec N$.
For incomparable Blocks sharing genesis, the permitted alternatives are given below.
\[
\begin{aligned}
\mathrm{Improves}(N,O)\iff{}&
 \bigl[\qid{2}(O)\prec\qid{2}(N)\bigr]\\
 &\lor\bigl[R(N,O)\\
 &\land\bigl(\qid{1}(O)\preceq\qid{1}(N)\\
 &\hspace{7mm}\lor\qid{2}(O)\preceq\qid{2}(N)\bigr)\bigr]\\
 &\lor\bigl[\qid{2}(O)\prec\qid{1}(N)\\
 &\hspace{7mm}\land\neg(\qid{2}(N)\prec\qid{1}(O))\bigr]\\
 &\lor\bigl[\qid{3}(O)\prec\qid{3}(N)\\
 &\hspace{5mm}\land\height(\qid{2}(O))\leq\height(\qid{3}(N))\bigr].
\end{aligned}
\]
The first alternative permits strict Q2 advance without a better score.
The score alternative requires Q1 or Q2 ancestry in addition to $R(N,O)$.
The negated condition retains equality, reverse ancestry and incomparable lineages as the exact complement of strict ancestry.

The final alternative permits an advancing Q3 that reaches the height of the old Q2 Block.
Its completed quorum support must satisfy the conditions in Section~\ref{sec:q3-support}.
Ancestry and score alternatives preserve Q2; this Q3 alternative can preserve finality while Q2 moves to an incomparable lineage.
The new Vote retains PrevSelfVote while its earlier linked history remains active.

Suppose children $B$ and $X$ share parent $A$, proposer Vote $p$ and a quorum of parent Votes.
Their scores are equal, and $\qid{1}(B)=\qid{1}(X)=A$.
If $B$ has the smaller identity, the score and Q1 conditions permit improvement from $X$ to $B$.
Votes on distinct Blocks sharing PrevVote follow the PrevSelfVote and active linked history checks for permitted signed order.

\section{Safety and Liveness}\label{sec:formal}

\subsection{Evidence and authority}\label{sec:evidence}\label{sec:authority}

Evidence identifies violations by a signing Validator through authenticated Votes and their references.
Direct evidence includes distinct Votes sharing PrevVote and the same exact endorsed Block, or distinct successors sharing a nonempty PrevSelfVote.
Other pairs require mutually active endorsed Blocks and neither Vote in the other's valid bounded linked history.
Derived evidence identifies a Vote violating required evidence or direction rules, using its signed references and verified supporting evidence.
A fresh Vote includes its required evidence and the available verified evidence.
Each Validator's stake is excluded once, even when several verified pieces of evidence identify that Validator.

A child Block records verified evidence and excludes the identified Validators from its authority.
For parent authority $s_P$ and identified Validator set $E_B$, the child's authority is $s_B=s_P\setminus E_B$.
Admission and parent Vote counting continue to use $s_P$; later children use $s_B$.
Historical authority remains unchanged, and $s_B$ must contain at least one Validator.

\subsection{Assumptions and Safety}\label{sec:assumptions}

\subsubsection{Finalized history}\label{sec:safety-result}

Safety requires compatible finalized histories across honest nodes, while preservation requires each node's later history to extend its earlier history.
These conditions make finalized Blocks stable references for later finalization.
Let $O(s,i)$ be honest node $i$'s finalized lineage at state $s$ of execution $H$.
Write $L\sqsubseteq R$ when $L$ is a prefix of $R$.
For all honest nodes $i,j$ and states $s,t$ in $H$, require the following relations.
\[
\begin{aligned}
\mathrm{Safety}:\quad
&O(s,i)\sqsubseteq O(t,j)\ \lor\ O(t,j)\sqsubseteq O(s,i),\\
\mathrm{Preservation}:\quad
&s\leq_H t\ \Longrightarrow O(s,i)\sqsubseteq O(t,i).
\end{aligned}
\]
The theorem establishes compatible finalized Blocks within each reachable state under the following assumptions.

\subsubsection{Quorum overlap}\label{sec:safety-quorums}

The Safety argument uses the honest Validator shared by any two quorums.
The improvement rule in Section~\ref{sec:permission} preserves Q2 ancestry or requires established Q3 support for catch-up.
Assume authenticated Votes, a common genesis and an initial authority with distinct Validators holding positive stake.
Honest Validators hold strictly more than two thirds of the initial stake and check improvement against every earlier Vote they recorded.
Verified exclusions preserve honest Validators, so quorums under independently reduced historical authorities derived from that initial authority share an honest Validator.
The theorem requires a quorum for child Block admission; Section~\ref{sec:proposals} defines the protocol's configurable endorsement threshold.

\subsubsection{Q3 support for catch-up}\label{sec:q3-support}

The Q3 catch-up rule in Section~\ref{sec:permission} uses history that honest Validators have already validated.
For finalized Block $F$, let $C$ be the Block at which its quorum support is established under the applicable historical authority.
Each contributing recorded Vote $v$ must retain $F$ in its Q2 ancestry and endorse a Block below $C$.
\[
F\preceq\qid{2}(T(v)),\qquad
\height(T(v))<\height(C).
\]
For non-genesis Q3 $F$, the node must already have admitted recorded Votes from a quorum, each satisfying both inequalities above.
Block $C$ must have height at most that of the new endorsed Block.
The height bounds tie the new Vote to earlier validated Votes, preserving finalized ancestry during catch-up.

\subsection{Conditional progress}\label{sec:liveness}

Conditional progress assumes that every honest node eventually receives the same finite ordered collection of proposed Block lineages.
Its first lineage is valid and precedes every other listed lineage under the Block comparison.
Required Blocks and Votes, including references retained by any honest node, must eventually reach every honest node.
Delivery permits arbitrary finite delays, duplicates and reordered messages.
This differs from partial synchrony, which adds timing bounds under the conditions studied by Dwork, Lynch and Stockmeyer~\citep{rw-dls}.

Honest Validators must hold strictly more than two thirds of the stake in every applicable parent authority.
Each honest Validator supplies a valid Vote in each sequence used to finalize the selected Block and a descendant.
These Votes have increasing recorded order, distinct exact endorsed Blocks per Validator, and permitted changes relative to earlier recorded Votes.
Each non-genesis Q3 has already admitted support with the height bound in Section~\ref{sec:q3-support}.
Each contributing Vote satisfies $F\preceq\qid{2}(T(v))$ and $\height(T(v))<\height(C)$ for the Block $F$ being finalized.
Fresh generation supplies a valid strict descendant with sufficient parent Votes and an authority derived from verified evidence.

Under these assumptions, the result constructs a finite extension of the same execution that finalizes the selected valid Block and a valid strict descendant.
Both finalized Blocks occur in one resulting state, advancing finalized history beyond the selected Block.
The extension permits arbitrarily long finite message delays.
Liveness requires continued finalization above every common reference height, as defined in Section~\ref{sec:overview}.

\section{Related research}\label{sec:related}

\subsection{Finalization}

Chained HotStuff uses a three-chain commit rule under partial synchrony~\citep[Sections 4--5]{rw-hotstuff}.
Its leader coordinates the Votes that establish successive quorum certificates.
Tsampi applies Q1 three times, restarting from each returned Block.
Each query accumulates distinct Validators' stake along that Block's ancestry under the authority recorded there.

\subsection{Permitted Votes}

Tendermint orders proposals through rounds with prevote and precommit steps~\citep[Section III, Algorithm 1]{rw-tendermint}.
Its locks constrain which proposal a validator may support in a later round.
Tsampi checks PrevSelfVote and the ancestry of historical quorum results before permitting a later Vote.
PrevSelfVote identifies the earlier signed choice against which the improvement rule is checked.

Casper FFG defines slashing conditions for Votes between source and target checkpoints~\citep[Section 2]{rw-casper}.
Its accountable Safety theorem connects conflicting finalized checkpoints to violations by at least one third of the total deposit.
Tsampi checks conflicts between authenticated Votes and applies verified evidence to a child Block's authority.
Tsampi uses the parent authority for admission, preserving historical Vote weights while verified evidence determines the child's authority.

\subsection{Delivery}

Narwhal separates transaction dissemination from ordering, and Tusk provides asynchronous consensus over the resulting directed acyclic graph~\citep[Sections 3 and 5]{rw-narwhal-tusk}.
Narwhal's certificates of availability establish that honest participants retain the referenced transaction data.
Tsampi describes single-parent Block lineages and assumes eventual delivery of required Blocks and Votes.
Historical quorum queries require the parent Votes and exact Block ancestry, so delivery includes those references.

\begingroup\raggedright
\bibliography{references}
\endgroup
\end{document}
